\scanpage{128}
\mspage{63}
Il faut introduire encore les espaces analytiques
de drapeaux $D_{d_0,d_1,\ldots,d_r}$ pour $d_0\geq d_1\geq\cdots\geq d_r\geq0$
leur somme est notée $D_{\Delta_r}$, où $\Delta_r$ désigne
maintenant l'ens. ordonné
\[
\tag{203}\Delta_r=[0,r]\cap\mathbf N
\]
ou simplexe ordonné type de dim $r$. Cette notation
$D_{\Delta_r}$ est compatible avec la notation $D_J$ pour $J$
ens. fini ordonné ayant un plus petit élément.
Alors $D_{\unsure{I_r|\Delta_r}}$ dépend fonctoriellement de façon
contravariante en $\unsure{I_r|\Delta_r}$, pour les applications
croissantes, entre simplexes, donc on a un
espace analytique $1/2$-simplicial, qu'on désigne par
$\mathcal D_*$, donc
\[
\tag{204}
\underset{\inserted{\text{ou }D_{\Delta_r}}}{\mathcal D_r}
=\deleted{\unsure{D_{\Delta_I}|D_{\Delta_r}}}\inserted{D_{\Delta_r}}
=\coprod_{\substack{(d_0,\ldots,d_r)\in\mathbf N^r\\d_0>d_1\cdots>d_r}}
D_{d_0,d_1,\ldots,d_r}.
\]
\deleted{$\unsure{D_{d_0,d_r}|D_{\Delta_r,d_r}}$}
La composante $D_{d_*}=D_{d_0,d_1,\ldots,d_r}$ est
un revêtement étale fini de $D_{d_0}$, s'identifiant
à
\[
\tag{205}
D_{d_*}=D'_{d_0}\wedge_{(\unsure{\mathfrak S_{d_0}|G_{d_0}})_{D_{d_0}}}
\underbrace{\mathcal P_{d_*}\bigl(\deleted{\unsure{[1,d_0]}}I_{d_0}\bigr)}_{
\substack{\text{ens des « drapeaux de}\\
\text{type }d_*\text{ » de l'ens fini}\\I_{d_0}=[1,d_0]\cap\mathbf N}}
\]
on désigne par $D^*_{d_0,d_1,\ldots,d_r}$ l'image inverse
de $D^*_{d_0}$, qui est donc le complémentaire
de l'image inverse $\Delta_{d_0- d_r}$ du diviseur
$\Theta_{d_0}$ sur $D_{d_0}$. On peut bien sûr plus \unsure{simple-|simple}
\margin{127}
